\section{Evaluation}
\label{sec:eval}







\subsection{Comparison with State-of-the-Art}
\label{subsec:main_results}



\begin{table*}[tb]
    \centering
    \caption{\textbf{Instruction-following capability comparison} with state-of-the-art text-to-motion models. Motion categories: (a) Locomotion, (b) Sports \& Athletics, (c) Fitness \& Outdoor Activities, (d) Daily Activities, (e) Social Interactions \& Leisure, and (f) Game Character Actions.}
    \vspace{5pt}
    \label{tab:sota_comparison}
    \begin{tabular}{l|cccccc|c|c}
        \toprule
        Method & Catg. (a) & (b) & (c) & (d) & (e) & (f) & Avg. & SSAE  \\
        \midrule
        MoMask \cite{momask} & 2.98 & 2.41 & 2.09 & 2.07 & 2.38 & 1.97 & 2.31 & 58.0\% \\
        GoToZero \cite{Fan_2025_ICCV} & 2.80 & 2.23 & 2.07 & 2.00 & 2.32 & 1.74 & 2.19 & 52.7\% \\
        DART \cite{DartControl} & 2.91 & 2.47 & 2.03 & 2.07 & 2.40 & 2.05 & 2.31  & 42.7\% \\
        LoM \cite{chen2024language} & 2.81 & 2.07 & 1.95 & 2.00 & 2.39 & 1.84 & 2.17  & 48.9\% \\
        \midrule
        \textbf{\modelname{}} & \textbf{3.76} & \textbf{3.18} & \textbf{3.15} & \textbf{3.06} & \textbf{3.25} & \textbf{3.01} & \textbf{3.24} & \textbf{78.6\%} \\
        \bottomrule
    \end{tabular}
    \vspace{-10pt}
\end{table*}


\begin{table*}[tb]
    \centering
    \caption{\textbf{Motion quality comparison} with state-of-the-art text-to-motion models. Motion categories: (a) Locomotion, (b) Sports \& Athletics, (c) Fitness \& Outdoor Activities, (d) Daily Activities, (e) Social Interactions \& Leisure, and (f) Game Character Actions.}
    \vspace{5pt}
    \label{tab:sota_comparison_mq}
    \begin{tabular}{l|cccccc|c}
        \toprule
        Method & Catg. (a) & (b) & (c) & (d) & (e) & (f) & Avg.  \\
        \midrule
        MoMask \cite{momask} & 3.05 & 2.91 & 2.58 & 2.66 & 2.77 & 2.81 & 2.79 \\
        GoToZero \cite{Fan_2025_ICCV} & 3.11 & 3.01 & 2.69 & 2.72 & 2.89 & 2.81 & 2.86  \\
        DART \cite{DartControl} & 3.38 & 3.33 & 2.94 & 2.95 & 3.06 & 3.07 & 3.11  \\
        LoM \cite{chen2024language} & 3.14 & 3.08 & 2.98 & 3.01 & 3.14 & 3.01 & 3.06  \\
        \midrule
        \textbf{\modelname{}} & \textbf{3.59} & \textbf{3.51} & \textbf{3.28} & \textbf{3.37} & \textbf{3.43} & \textbf{3.41} & \textbf{3.43} \\
        \bottomrule
    \end{tabular}
\end{table*}


To evaluate \modelname{} against state-of-the-art text-to-motion models, we constructed a diverse test set comprising over 2000 text prompts. This test set spans six major categories (\cref{subsec:stats}) and covers simple atomic actions as well as concurrent and sequential action combinations. Human annotators were asked to rate the generated motions on a scale of 1 to 5 from two perspectives: instruction-following capability and motion quality. Fig. \ref{fig:teasercomp} and Tabs. \ref{tab:sota_comparison}–\ref{tab:sota_comparison_mq} show the comparison with four models: DART \cite{DartControl}, LoM \cite{chen2024language}, GoToZero \cite{Fan_2025_ICCV}, and MoMask \cite{momask}. Fig. \cref{fig:qualitative} presents some examples of the visual comparison. Our model significantly outperforms the other models in terms of both instruction-following capability and motion quality.

We further adopt an automatic evaluation approach, namely Structured Semantic Alignment Evaluation (SSAE) \cite{cao2025hunyuanimage}, which utilizes powerful video-VLMs to assess the generated motions. This approach transforms the text-motion alignment problem into a video-question-answering task. It works by decomposing each input prompt into a series of ``yes'' or ``no'' questions. These questions are then presented to a video-VLM (\eg, Gemini-2.5-Pro) along with the rendered video of the generated motion. For example, given the prompt ``a person swings their arm while shooting a soccer ball'', the decomposed questions might include: ``is the person kicking their leg?'', ``is the person swinging their arm?'', and ``does the person appear to be shooting a soccer ball?''. The correct rate of the VLM's answers across all test prompts constitutes the model's SSAE score. A comparison of SSAE scores with other models is provided in \cref{tab:sota_comparison}.



\begin{figure}[tb]
    \centering
    \begin{tabular}{c}
        \includegraphics[width=0.9\linewidth, trim=50 450 45 30]{figures/exp/comparison_page1.pdf} \\ \\
        \includegraphics[width=0.9\linewidth, trim=50 450 45 30]{figures/exp/comparison_page2.pdf}
    \end{tabular}
    \vspace{10pt}
    \caption{
        Examples of visual comparison to state-of-the-art models.
    }
    \label{fig:qualitative}
\end{figure}



\subsection{Scaling Experiments}
\label{subsec:ablation}

To understand the effects of scale, we investigated how model size and data volume impact performance. The results of our human evaluation on models of varying sizes are presented in Tabs. \ref{tab:scalingexp}–\ref{tab:scalingexp_mq}. Our findings indicate that while instruction-following capability consistently improves with larger models, motion quality reaches a saturation point beyond the 0.46B parameter size. Notably, a comparison between the "DiT-0.46B" and "DiT-0.46B-400h" models (Tab. \ref{tab:scalingexp}) underscores the critical role of the larger 3,000-hour training dataset in enhancing instruction-following capability. Consequently, we plan to release "DiT-1B" as the primary HY-Motion 1.0 model and "DiT-0.46B" as a lightweight "Lite" version.


\begin{table*}[tb]
    \centering
    \caption{\textbf{Instruction-following capability comparison} of different model sizes. Motion categories: (a) Locomotion, (b) Sports \& Athletics, (c) Fitness \& Outdoor Activities, (d) Daily Activities, (e) Social Interactions \& Leisure, and (f) Game Character Actions. The model ``DiT-0.46B-400h'' is trained only on the 400-hour high-quality dataset, while the other models are pretrained on the 3,000-hour dataset.}
    \vspace{5pt}
    \label{tab:scalingexp}
    \begin{tabular}{l|cccccc|c}
        \toprule
        Method & Catg. (a) & (b) & (c) & (d) & (e) & (f) & Avg.   \\
        \midrule
        DiT-0.05B & 3.68 & 2.95 & 2.85 & 3.10 & 3.12 & 2.90 & 3.10  \\
        DiT-0.46B & 3.93 & 3.15 & 3.02 & 3.10 & 3.15 & 2.88 & 3.20 \\
        DiT-0.46B-400h & 3.76 & 2.95 & 2.73 & 2.93 & 3.17 & 2.75 & 3.05 \\
        DiT-1B & \textbf{3.95} & \textbf{3.32} & \textbf{3.23} & \textbf{3.36} & \textbf{3.27} & \textbf{2.92} & \textbf{3.34}  \\
        \bottomrule
    \end{tabular}
    \vspace{-10pt}
\end{table*}


\begin{table*}[tb]
    \centering
    \caption{\textbf{Motion quality comparison} of different model sizes. Motion categories: (a) Locomotion, (b) Sports \& Athletics, (c) Fitness \& Outdoor Activities, (d) Daily Activities, (e) Social Interactions \& Leisure, and (f) Game Character Actions.}
    \vspace{5pt}
    \label{tab:scalingexp_mq}
    \begin{tabular}{l|cccccc|c}
        \toprule
        Method & Catg. (a) & (b) & (c) & (d) & (e) & (f) & Avg.  \\
        \midrule
        DiT-0.05B & 3.32 & 2.54 & 2.95 & 2.86 & 3.00 & 2.80 & 2.91  \\
        DiT-0.46B & 3.68 & \textbf{3.32} & 3.10 & 3.10 & 3.22 & 3.13 & 3.26 \\
        DiT-0.46B-400h & \textbf{3.83} & 3.02 & 3.23 & \textbf{3.29} & 3.27 & 3.20 & 3.31 \\
        DiT-1B & 3.78 & 3.20 & \textbf{3.33} & 3.26 & \textbf{3.27} & \textbf{3.23} & \textbf{3.34}  \\
        \bottomrule
    \end{tabular}
\end{table*}










